\documentclass[a4paper,11pt]{article}

\usepackage[margin=22mm]{geometry}
\usepackage{amsmath,amssymb}
\usepackage{booktabs}
\usepackage{array}
\usepackage{tabularx}
\usepackage{hyperref}
\usepackage{zxjatype}
\usepackage[haranoaji]{zxjafont}

\setlength{\parindent}{0pt}
\setlength{\parskip}{0.45em}
\renewcommand{\arraystretch}{1.15}
\newcolumntype{L}[1]{>{\raggedright\arraybackslash}p{#1}}
\newcommand{\reraptor}{ReRAPTOR}

\begin{document}
\sloppy

\title{式(3)および4.2節の記号・定数に関する補足説明}
\section*{1. 式(3)の意味と $c_i,\alpha_i$ の設定理由}

式(3)の正規化付き geometry-aware 既定インスタンスは
\[
F_{\mathrm{geo}}(r)=\sum_{j=1}^{4}\alpha_j \tilde{f}_j(r),
\qquad
\tilde{f}_j(r)=\frac{f_j(r)}{c_j}
\]
であり，ここで $c_i$ は各評価項目のスケールを揃えるための正規化定数です。

また，本文で明示している既定の評価重みは
\[
(\alpha_1,\alpha_2,\alpha_3,\alpha_4)=(0.40,0.25,0.20,0.15)
\]
であり，それぞれ ride duration，unique visited stops，convex-hull area，quadrant coverage に対応します。

\begin{tabularx}{\linewidth}{L{2.2cm}L{2.3cm}L{3.8cm}X}
\toprule
項目 & 値 & 意味 & 設定理由・由来 \\
\midrule
$\alpha_1$ & $0.40$ & ride duration の重み & geometry-aware 既定プロファイルでは，まず「十分な乗車体験」を主成分に置くため最大に設定している。福岡未踏システムで用いていた balanced profile を継承した経験的既定値であり，本稿で再最適化した値ではない。 \\
$\alpha_2$ & $0.25$ & unique visited stops の重み & Stop 多様性を第2の価値として評価するための重みである。ride duration を主成分に据えつつ，多様性を無視しないための既定比率であり，初期実装由来の経験的既定値である。 \\
$\alpha_3$ & $0.20$ & convex-hull area の重み & 空間的な広がりを補助的に評価するための重みである。ride と stop diversity よりは低く置き，geometry-aware な既定インスタンスの補助指標として扱っている。初期実装由来の経験的既定値である。 \\
$\alpha_4$ & $0.15$ & quadrant coverage の重み & 四方向への広がりを補助的に評価するための重みである。$f_4$ は二値的・粗い指標なので，$f_3$ よりさらに控えめに置いている。初期実装由来の経験的既定値である。 \\
$c_1$ & $480$ & ride duration の正規化定数 & 最大予算 8 時間（480 分）に対応する物理的上限であり，問題設定から直接決まる。 \\
$c_2$ & $50$ & unique visited stops の正規化定数 & Stop 数を他項目と同程度のスケールに揃えるための参照値である。物理定数ではなく，初期実装で置いた正規化用参照値であり，本稿で最適化した値ではない。 \\
$c_3$ & $100$ & convex-hull area の正規化定数 & hull area の量的スケールが他項目を過度に支配しないように置いた参照値である。これも正規化用参照値であり，最適化結果ではない。 \\
$c_4$ & $4$ & quadrant coverage の正規化定数 & 四象限カバレッジの理論上限であり，問題設定から直接決まる。 \\
\bottomrule
\end{tabularx}

したがって，$c_1,c_4$ は問題設定由来，$c_2,c_3$ は正規化参照値，
$\alpha_1,\ldots,\alpha_4$ は geometry-aware 既定プロファイルの重みであり，
いずれも「本稿の比較実験で最適化して得た値」としては扱っていません。

\section*{2. 4.2節の記号・関数の説明}

\subsection*{2.1 ラベル $\lambda$ の各要素}

\[
\lambda = \left(
a,\rho,d,V,q,\Pi,c^{\mathrm{stop}},c^{\mathrm{line}},n^{\mathrm{tr}},\psi
\right)
\]

\begin{tabularx}{\linewidth}{L{2.2cm}L{4.3cm}X}
\toprule
記号 & 意味 & 用途 \\
\midrule
$a$ & 現在の到着時刻 & 次に乗車可能な便の判定 \\
$\rho$ & 累積乗車時間（分） & 乗車価値の評価 \\
$d$ & 累積移動距離（km） & 距離指標として保持 \\
$V$ & 訪問済み停留所集合 & Stop 数や空間被覆の計算 \\
$q$ & 象限カバレッジのビットマスク & 東西南北 4 象限の訪問状態 \\
$\Pi$ & 区間列（乗車履歴） & 最終経路の復元 \\
$c^{\mathrm{stop}}$ & Stop ごとの訪問回数カウンタ & 過訪検出，Negative 枝刈り \\
$c^{\mathrm{line}}$ & 路線ごとの乗車回数カウンタ & 路線テーマ制約の補助 \\
$n^{\mathrm{tr}}$ & 乗換回数 & 経路複雑度の補助指標 \\
$\psi$ & 目的固有アキュムレータ & stop-reward タスクでは獲得報酬総和 $\sigma$ と獲得 Stop 数 $m$ \\
\bottomrule
\end{tabularx}

注記として，実装上は $\psi$ を抽象変数としてそのまま保持せず，
`objective\_score\_total` と `objective\_scored\_count` に分けて保持しています。

\subsection*{2.2 拡張演算子の各要素}

時刻表上の 1 区間を
\[
e=(u,v,\tau^{\mathrm{dep}},\tau^{\mathrm{arr}},\ell,d_e)
\]
とします。ここで，
\begin{itemize}
    \item $u, v$ は出発 Stop と到着 Stop
    \item $\tau^{\mathrm{dep}}, \tau^{\mathrm{arr}}$ は発車時刻と到着時刻
    \item $\ell$ は路線 ID
    \item $d_e$ は区間距離
\end{itemize}
です。

また，
\begin{itemize}
    \item $q(v)$ は Stop $v$ が属する象限を表すビット
    \item $\Pi \oplus e$ は区間列 $\Pi$ に区間 $e$ を追加する操作
    \item $U_{\theta}(\psi,\lambda,e)$ は目的固有の状態更新関数
\end{itemize}
です。

stop-reward タスクでは，$U_{\theta}$ は新しい Stop に初めて到達したときだけ
対応する報酬を加算します。

\subsection*{2.3 4.2節で出てくる主な関数}

\begin{tabularx}{\linewidth}{L{3.0cm}X}
\toprule
記号 & 意味 \\
\midrule
$F_{\theta}(r)$ & シナリオ $\theta$ における経路価値関数 \\
$g_{\theta}(r)$ & シナリオ $\theta$ における受理条件 \\
$A(V)$ & 訪問 Stop 座標から得た凸包面積 \\
$Q(V)$ & カバーした象限数 \\
$S_{\theta}(\lambda)$ & 探索中のラベル評価スコア \\
$P(\lambda)$ & Negative penalty \\
$P_{\theta}(\lambda)$ & タスク依存 penalty \\
$D_{\theta}(\lambda^a,\lambda^b)$ & 支配性判定関数 \\
$\operatorname{TopW}$ & ビーム幅 $W$ で上位ラベルのみ残す演算 \\
$\operatorname{Next}(u,a)$ & Stop $u$ から時刻 $a$ 以降で利用可能な候補区間集合 \\
\bottomrule
\end{tabularx}

\section*{3. 4.2節に現れる数値の意味と設定理由}

本節では，4.2節に現れる数値について，
「何を制御している値か」「なぜそのように置いたか」「経験か実験結果か」を
本文に対応する形で整理します。

\subsection*{3.1 Geometry-aware 探索スコアの係数}

geometry-aware な探索スコアは
\[
S_{\mathrm{geo}}(\lambda) =
\beta_1 \rho(\lambda)
+ \beta_2 |V(\lambda)|
+ \beta_3 A(V(\lambda))
+ \beta_4 Q(V(\lambda))
- P(\lambda)
\]
であり，既定係数は
\[
(\beta_1,\beta_2,\beta_3,\beta_4)=(1.5,10,50,100)
\]
です。

\begin{tabularx}{\linewidth}{L{2.0cm}L{2.0cm}L{3.5cm}X}
\toprule
項目 & 値 & 役割 & 設定理由・由来 \\
\midrule
$\beta_1$ & $1.5$ & ride time の加点係数 & 乗車時間を正の価値として扱う基本係数である。ただし geometry-aware 既定インスタンスでは ride time 単独よりも空間拡張を強く見たいので，他係数より小さく置いている。初期実装由来の経験的既定値。 \\
$\beta_2$ & $10$ & 訪問 Stop 数の加点係数 & 新規 Stop を増やす方向を ride time より強く促すための係数である。局所的な往復より Stop 多様性を優先させるため，$\beta_1$ より大きく置いている。初期実装由来の経験的既定値。 \\
$\beta_3$ & $50$ & hull area の加点係数 & 面的な広がりを強く促すための係数である。geometry-aware 既定インスタンスでは，局所的に Stop 数だけを増やすより空間被覆を優先したいため，大きく置いている。初期実装由来の経験的既定値。 \\
$\beta_4$ & $100$ & quadrant coverage の加点係数 & 探索初期に一方向へ偏ることを避け，まず複数方向へ広がる partial route を優先するため最大に設定している。粗い指標だが探索の初期誘導には有効という設計意図に基づく経験的既定値。 \\
\bottomrule
\end{tabularx}

これらは本稿で網羅的に調整した最適値ではなく，
geometry-aware 既定インスタンスの探索誘導を安定させるために
初期実装で採用していた経験的既定値です。

\subsection*{3.2 Stop-reward 探索スコアの係数}

Exp1--3 で用いている stop-reward 側の探索スコアは
\[
\begin{aligned}
S_{\omega,\eta}(\lambda)
=\;& 10000 \, \sigma(\lambda)
+ 5 \, \operatorname{SPH}(\lambda)
- 0.1 \, T_{\eta}(\lambda) \\
&- 50000 \, \max\!\left(0, T_{\eta}(\lambda)-B\right)
- P(\lambda),
\end{aligned}
\]
であり，
\[
\operatorname{SPH}(\lambda)=60\,\sigma(\lambda)/T_{\eta}(\lambda)
\]
です。

\begin{tabularx}{\linewidth}{L{2.1cm}L{2.0cm}L{3.4cm}X}
\toprule
項目 & 値 & 役割 & 設定理由・由来 \\
\midrule
$10000$ & $10000$ & $\sigma(\lambda)$ の係数 & collected reward を主目的として最優先させるための大きな係数である。小さな reward 差が時間系の補助項より常に優先されるよう，辞書式優先順位をスカラー値で近似している。現在の stop-reward 課題の設計上の既定値であり，最適化結果ではない。 \\
$5$ & $5$ & $\operatorname{SPH}(\lambda)$ の係数 & 同じ reward 総量なら，単位時間当たりの reward 効率が高い候補を優先するための tie-breaker 係数である。reward を主目的に残したまま，効率だけを二次基準として使うため控えめに置いている。現在の課題設計上の既定値。 \\
$0.1$ & $0.1$ & $T_{\eta}(\lambda)$ の減点係数 & reward と効率が同程度のときに，より短い effective time の partial route をわずかに優先するための微弱な減衰係数である。主目的を壊さないよう小さく置いている。現在の課題設計上の既定値。 \\
$50000$ & $50000$ & budget 超過ペナルティ係数 & $T_{\eta}(\lambda)>B$ の候補を事実上排除するための大きな係数である。reward の正項より十分大きくし，budget violation を強く抑止するために用いている。現在の課題設計上の既定値であり，実験による最適化値ではない。 \\
$60$ & $60$ & $\operatorname{SPH}$ の換算係数 & 分単位の effective time から「1時間あたりの score」に直す単位変換係数である。これはチューニング値ではなく，単位変換のための定数である。 \\
\bottomrule
\end{tabularx}

ここで使っている数値は，reward $\gg$ efficiency $\gg$ time という
優先順位を実装するための設計上の既定値であり，
比較実験から逆算した最適パラメータではありません。

\subsection*{3.3 Negative penalty の閾値と係数}

Negative penalty は
\[
P(\lambda)=
50000\,\mathbf{1}\!\left[\max_{s} c^{\mathrm{stop}}_{\lambda}(s) > 2\right]
+ 20000\,\mathbf{1}\!\left[\rho(\lambda) / |V(\lambda)| < 2\right]
+ P_{\theta}(\lambda)
\]
です。

\begin{tabularx}{\linewidth}{L{2.2cm}L{2.0cm}L{3.5cm}X}
\toprule
項目 & 値 & 役割 & 設定理由・由来 \\
\midrule
$>2$ & $2$ を超える & 同一 Stop の過訪検出 & 同じ Stop を 3 回以上含む partial route は，不自然な巡回に入り始めている可能性が高いので排除対象とみなす閾値である。初期実装から継承した経験的閾値であり，本稿の最適化結果ではない。 \\
$50000$ & $50000$ & 過訪ラベルへの強ペナルティ & 上記の過訪条件を満たしたラベルを，通常の加点項より強く抑えるための支配的ペナルティ係数である。明らかな loop を確実に落とすための安全側の既定値。 \\
$<2$ & $2$ 未満 & 極端に短い hop の連続検出 & 1 Stop あたり平均乗車時間が 2 分未満という条件は，短い hop を連続して積み重ねる不自然な経路を検出するための閾値である。これも経験的既定値であり，最適化した閾値ではない。 \\
$20000$ & $20000$ & 短 hop 連続ラベルへの強ペナルティ & 上記条件を満たしたラベルを抑えるための支配的ペナルティ係数である。$50000$ よりは弱いが，依然として通常の加点項より十分大きく置いている。初期実装由来の経験的既定値。 \\
$P_{\theta}(\lambda)$ & タスク依存 & シナリオ依存 penalty & theme や scenario 制約を入れる場合だけ追加する penalty であり，値そのものはシナリオ側で決まる。したがって固定の最適化値ではなく，課題依存の追加項である。 \\
\bottomrule
\end{tabularx}

\subsection*{3.4 支配性判定に出てくる閾値}

geometry-aware 側の支配性判定は
\[
\lambda^{a} \succeq_{\mathrm{geo}} \lambda^{b}
\iff
S_{\mathrm{geo}}(\lambda^{a}) \ge S_{\mathrm{geo}}(\lambda^{b})
\;\land\;
\rho(\lambda^{a}) \ge 0.95\,\rho(\lambda^{b})
\;\land\;
A(V(\lambda^{a})) \ge A(V(\lambda^{b})).
\]
であり，ここで唯一の数値閾値は $0.95$ です。

\begin{tabularx}{\linewidth}{L{2.0cm}L{2.2cm}L{3.6cm}X}
\toprule
項目 & 値 & 役割 & 設定理由・由来 \\
\midrule
$0.95$ & $0.95$ & ride time 近似保持の許容率 & geometry-aware な支配性では，ride time がわずかに短いだけで area や total score に優れるラベルまで落とし過ぎないよう，完全一致ではなく 95\% 以上を許容する緩和閾値である。探索を過度に硬くしないための経験的既定値であり，本稿で最適化した値ではない。 \\
\bottomrule
\end{tabularx}

\subsection*{3.5 Beam 幅}

beam pruning は
\[
\mathcal{L}_{k+1}(v) =
\operatorname{TopW}_{S_{\theta}}
\Big(
\{\lambda' \mid \lambda \in \mathcal{L}_{k}(u),\;
e \in \operatorname{Next}(u,a_{\lambda}),\;
\lambda'=\operatorname{Extend}(\lambda,e)\}
\Big)
\]
であり，本文では既定値として
\[
W=50
\]
を用いています。

\begin{tabularx}{\linewidth}{L{2.0cm}L{2.2cm}L{3.5cm}X}
\toprule
項目 & 値 & 役割 & 設定理由・由来 \\
\midrule
$W$ & $50$ & 各 stop-round bucket に残すラベル数 & 探索品質と計算量の極端な偏りを避けるための固定既定値である。beam を狭くし過ぎると探索が早期に偏り，広くし過ぎると計算量が急増するため，中庸の既定値として用いている。初期実装から継承した経験的既定値であり，本稿で最適化した値ではない。 \\
\bottomrule
\end{tabularx}

\section*{4. 4.2節の stop-reward 支配性判定の意味}

stop-reward 版の支配性判定は
\[
\lambda^{a} \succeq_{\omega,\eta} \lambda^{b}
\iff
\sigma(\lambda^{a}) \ge \sigma(\lambda^{b})
\;\land\;
T_{\eta}(\lambda^{a}) \le T_{\eta}(\lambda^{b})
\;\land\;
U_{\omega,\eta}(\lambda^{a}) \ge U_{\omega,\eta}(\lambda^{b}),
\]
であり，少なくとも 1 つが狭義不等式です。

\begin{tabularx}{\linewidth}{L{2.7cm}L{4.0cm}X}
\toprule
記号 & 意味 & このように定義する理由 \\
\midrule
$\sigma(\lambda)$ & 現時点までに獲得した stop reward 総和 & stop-reward 課題の主目的そのものなので，支配性判定でも最優先の比較軸とする。 \\
$T_{\eta}(\lambda)$ & dwell time を含む effective time & 同じ reward なら，より少ない有効時間で到達しているラベルの方が将来の拡張余地が大きいため，比較軸に入れる。 \\
$U_{\omega,\eta}(\lambda)$ & 残り有効時間内に取り得る報酬の楽観的上界 & 現在の reward が同程度でも，将来取り得る最大 reward が明らかに低いラベルは落とせるため，上界を入れている。 \\
少なくとも1つが狭義不等式 & 完全同値ラベルだけを機械的に消すのではなく，少なくともどこかで厳密に優れている場合だけ支配とみなす条件 & 同値なラベルまで不必要に落とさないための標準的な支配条件である。これはパラメータ最適化の結果ではなく，支配関係の定義である。 \\
\bottomrule
\end{tabularx}

\end{document}
